% Options for packages loaded elsewhere
% Options for packages loaded elsewhere
\PassOptionsToPackage{unicode}{hyperref}
\PassOptionsToPackage{hyphens}{url}
\PassOptionsToPackage{dvipsnames,svgnames,x11names}{xcolor}
%
\documentclass[
  indonesian,
  11pt,
  a4paper,
]{report}
\usepackage{xcolor}
\usepackage[top=2.2cm,bottom=2.2cm,left=2.5cm,right=2.5cm]{geometry}
\usepackage{amsmath,amssymb}
\setcounter{secnumdepth}{5}
\usepackage{iftex}
\ifPDFTeX
  \usepackage[T1]{fontenc}
  \usepackage[utf8]{inputenc}
  \usepackage{textcomp} % provide euro and other symbols
\else % if luatex or xetex
  \usepackage{unicode-math} % this also loads fontspec
  \defaultfontfeatures{Scale=MatchLowercase}
  \defaultfontfeatures[\rmfamily]{Ligatures=TeX,Scale=1}
\fi
\usepackage{lmodern}
\ifPDFTeX\else
  % xetex/luatex font selection
\fi
% Use upquote if available, for straight quotes in verbatim environments
\IfFileExists{upquote.sty}{\usepackage{upquote}}{}
\IfFileExists{microtype.sty}{% use microtype if available
  \usepackage[]{microtype}
  \UseMicrotypeSet[protrusion]{basicmath} % disable protrusion for tt fonts
}{}
\usepackage{setspace}
\makeatletter
\@ifundefined{KOMAClassName}{% if non-KOMA class
  \IfFileExists{parskip.sty}{%
    \usepackage{parskip}
  }{% else
    \setlength{\parindent}{0pt}
    \setlength{\parskip}{6pt plus 2pt minus 1pt}}
}{% if KOMA class
  \KOMAoptions{parskip=half}}
\makeatother
% Make \paragraph and \subparagraph free-standing
\makeatletter
\ifx\paragraph\undefined\else
  \let\oldparagraph\paragraph
  \renewcommand{\paragraph}{
    \@ifstar
      \xxxParagraphStar
      \xxxParagraphNoStar
  }
  \newcommand{\xxxParagraphStar}[1]{\oldparagraph*{#1}\mbox{}}
  \newcommand{\xxxParagraphNoStar}[1]{\oldparagraph{#1}\mbox{}}
\fi
\ifx\subparagraph\undefined\else
  \let\oldsubparagraph\subparagraph
  \renewcommand{\subparagraph}{
    \@ifstar
      \xxxSubParagraphStar
      \xxxSubParagraphNoStar
  }
  \newcommand{\xxxSubParagraphStar}[1]{\oldsubparagraph*{#1}\mbox{}}
  \newcommand{\xxxSubParagraphNoStar}[1]{\oldsubparagraph{#1}\mbox{}}
\fi
\makeatother


\usepackage{longtable,booktabs,array}
\newcounter{none} % for unnumbered tables
\usepackage{calc} % for calculating minipage widths
% Correct order of tables after \paragraph or \subparagraph
\usepackage{etoolbox}
\makeatletter
\patchcmd\longtable{\par}{\if@noskipsec\mbox{}\fi\par}{}{}
\makeatother
% Allow footnotes in longtable head/foot
\IfFileExists{footnotehyper.sty}{\usepackage{footnotehyper}}{\usepackage{footnote}}
\makesavenoteenv{longtable}
\usepackage{graphicx}
\makeatletter
\newsavebox\pandoc@box
\newcommand*\pandocbounded[1]{% scales image to fit in text height/width
  \sbox\pandoc@box{#1}%
  \Gscale@div\@tempa{\textheight}{\dimexpr\ht\pandoc@box+\dp\pandoc@box\relax}%
  \Gscale@div\@tempb{\linewidth}{\wd\pandoc@box}%
  \ifdim\@tempb\p@<\@tempa\p@\let\@tempa\@tempb\fi% select the smaller of both
  \ifdim\@tempa\p@<\p@\scalebox{\@tempa}{\usebox\pandoc@box}%
  \else\usebox{\pandoc@box}%
  \fi%
}
% Set default figure placement to htbp
\def\fps@figure{htbp}
\makeatother



\ifLuaTeX
\usepackage[bidi=basic,shorthands=off,provide=*]{babel}
\else
\usepackage[bidi=default,shorthands=off,provide=*]{babel}
\fi
\ifLuaTeX
  \usepackage{selnolig} % disable illegal ligatures
\fi


\setlength{\emergencystretch}{3em} % prevent overfull lines

\providecommand{\tightlist}{%
  \setlength{\itemsep}{0pt}\setlength{\parskip}{0pt}}



 


\makeatletter
\@ifpackageloaded{caption}{}{\usepackage{caption}}
\AtBeginDocument{%
\ifdefined\contentsname
  \renewcommand*\contentsname{Daftar Isi}
\else
  \newcommand\contentsname{Daftar Isi}
\fi
\ifdefined\listfigurename
  \renewcommand*\listfigurename{Daftar Gambar}
\else
  \newcommand\listfigurename{Daftar Gambar}
\fi
\ifdefined\listtablename
  \renewcommand*\listtablename{Daftar Tabel}
\else
  \newcommand\listtablename{Daftar Tabel}
\fi
\ifdefined\figurename
  \renewcommand*\figurename{Gambar}
\else
  \newcommand\figurename{Gambar}
\fi
\ifdefined\tablename
  \renewcommand*\tablename{Tabel}
\else
  \newcommand\tablename{Tabel}
\fi
}
\@ifpackageloaded{float}{}{\usepackage{float}}
\floatstyle{ruled}
\@ifundefined{c@chapter}{\newfloat{codelisting}{h}{lop}}{\newfloat{codelisting}{h}{lop}[chapter]}
\floatname{codelisting}{Daftar}
\newcommand*\listoflistings{\listof{codelisting}{Daftar Daftar}}
\makeatother
\makeatletter
\makeatother
\makeatletter
\@ifpackageloaded{caption}{}{\usepackage{caption}}
\@ifpackageloaded{subcaption}{}{\usepackage{subcaption}}
\makeatother
\usepackage{bookmark}
\IfFileExists{xurl.sty}{\usepackage{xurl}}{} % add URL line breaks if available
\urlstyle{same}
% fallback for those not using the hyperref driver hyperxmp:
\makeatletter
\@ifundefined{xmpquote}{\newcommand{\xmpquote}[1]{#1}}{}
\makeatother
\hypersetup{
  pdftitle={Policy Brief: Zonasi Intervensi Kebijakan UMKM Kota Mataram},
  pdflang={id},
  pdfkeywords={\xmpquote{UMKM}, \xmpquote{Kota
Mataram}, \xmpquote{Digitalisasi}, \xmpquote{Sertifikasi
Halal}, \xmpquote{Indeks Kesiapan Digital}, \xmpquote{Analisis
Spasial}, \xmpquote{Kebijakan Berbasis Data}},
  colorlinks=true,
  linkcolor={blue},
  filecolor={Maroon},
  citecolor={Blue},
  urlcolor={Blue},
  pdfcreator={LaTeX via pandoc}}


\title{Policy Brief: Zonasi Intervensi Kebijakan UMKM Kota Mataram}
\usepackage{etoolbox}
\makeatletter
\providecommand{\subtitle}[1]{% add subtitle to \maketitle
  \apptocmd{\@title}{\par {\large #1 \par}}{}{}
}
\makeatother
\subtitle{Mempercepat Digitalisasi dan Sertifikasi Halal melalui
Kebijakan Berbasis Data dan Wilayah}
\author{}
\date{}
\begin{document}
\maketitle


\setstretch{1.15}
\chapter*{Tim Peneliti}\label{tim-peneliti}

\begin{itemize}
\item
  \textbf{Indira Puteri Kinasih, Ph.D.}\\
  indiraputeri@uinmataram.ac.id
\item
  \textbf{Dr.~Any Tsalasatul Fitriyah}\\
  any.tsalasatul@uinmataram.ac.id
\item
  \textbf{Nurul Susianti, M.E.}\\
  nurulsusianti@uinmataram.ac.id
\item
  \textbf{Ulfatun Hasanah, M.PA}\\
  ulfatunhasanah@uinmataram.ac.id
\item
  \textbf{Nur Awalunnisa, M.SEI.}\\
  awalunnisa@uinmataram.ac.id
\end{itemize}

\newpage

\chapter*{Ringkasan}\label{ringkasan}

Usaha Mikro, Kecil, dan Menengah (UMKM) memiliki peran penting dalam
perekonomian Kota Mataram. Namun, penguatan UMKM tidak cukup dilakukan
melalui program pelatihan dan bantuan yang seragam untuk seluruh pelaku
usaha. Perbedaan kapasitas digital, pemanfaatan perdagangan elektronik,
kondisi infrastruktur, serta kesiapan dan hambatan sertifikasi halal
memerlukan intervensi yang lebih terarah.

Penelitian \emph{Zonasi Intervensi Kebijakan UMKM Kota Mataram: Pemetaan
Geospasial untuk Akselerasi Digitalisasi dan Sertifikasi Halal}
mengembangkan pendekatan berbasis data untuk mengidentifikasi
karakteristik UMKM, memetakan kesiapan digital, serta mengeksplorasi
hubungan antara kesiapan digital dan sejumlah faktor pendukung maupun
penghambat pada tingkat kelurahan. Analisis menggabungkan data
karakteristik UMKM, survei pelaku usaha, pemetaan geospasial, dan
pemodelan statistik dengan mempertimbangkan keterkaitan antarwilayah.

Hasil analisis menunjukkan adanya variasi kesiapan digital antarpelaku
usaha. Dalam model spasial yang digunakan, pemanfaatan perdagangan
elektronik memiliki hubungan positif dengan indeks kesiapan digital
UMKM. Kesiapan terkait halal juga menunjukkan estimasi hubungan positif,
sedangkan hambatan terkait halal menunjukkan arah hubungan negatif.
Estimasi pengaruh infrastruktur digital mengarah positif, tetapi
ketidakpastian estimasinya masih cukup besar sehingga hasilnya belum
memberikan bukti yang kuat mengenai hubungan tersebut. Temuan ini perlu
dipahami sebagai hubungan statistik, bukan bukti sebab-akibat.

Berdasarkan hasil tersebut, policy brief ini merekomendasikan lima arah
kebijakan: (1) menetapkan zonasi prioritas intervensi berdasarkan
karakteristik dan kebutuhan UMKM; (2) mengintegrasikan pelatihan digital
dengan pendampingan penggunaan perdagangan elektronik; (3) memperkuat
layanan konsultasi dan fasilitasi sertifikasi halal; (4) menghubungkan
pendampingan digital dan halal dengan penguatan mutu produk; serta (5)
membangun sistem pemantauan program berbasis data geospasial.

Pemerintah Kota Mataram bersama perangkat daerah terkait, pendamping
UMKM, perguruan tinggi, dan lembaga layanan halal perlu mengembangkan
pendekatan tersebut secara bertahap. Tujuannya bukan sekadar
meningkatkan jumlah kegiatan pelatihan atau sertifikat yang diterbitkan,
melainkan memperkuat kemampuan UMKM untuk memanfaatkan teknologi,
memenuhi persyaratan halal, dan mengembangkan usaha secara
berkelanjutan.

\chapter*{Latar Belakang Masalah}\label{latar-belakang-masalah}

Digitalisasi membuka peluang bagi UMKM untuk memperluas jangkauan pasar,
memperbaiki pengelolaan usaha, dan meningkatkan akses terhadap informasi
serta layanan ekonomi. Pada saat yang sama, pemenuhan persyaratan halal
menjadi bagian penting dalam pengembangan usaha, khususnya bagi pelaku
usaha makanan dan minuman yang perlu memperhatikan kepatuhan terhadap
ketentuan yang berlaku.

Meskipun demikian, kemampuan UMKM untuk memanfaatkan peluang tersebut
tidak selalu sama. Sebagian pelaku usaha mungkin telah menggunakan media
sosial atau platform perdagangan elektronik untuk menjual produk,
sementara sebagian lainnya masih memerlukan dukungan dalam keterampilan
digital, strategi pemasaran, pengelolaan informasi, dan pemanfaatan
teknologi secara konsisten.

Demikian pula, pengetahuan mengenai sertifikasi halal tidak selalu
diikuti oleh kesiapan untuk memenuhi persyaratan administratif maupun
teknis. Pelaku usaha dapat memerlukan penjelasan mengenai alur
sertifikasi, dokumen yang dibutuhkan, proses pendampingan, dan penerapan
jaminan produk halal dalam kegiatan produksi sehari-hari.

Pendekatan kebijakan yang seragam berisiko menghasilkan intervensi yang
kurang sesuai dengan kebutuhan aktual. Pelatihan pemasaran digital,
misalnya, belum tentu menjadi kebutuhan paling mendesak bagi UMKM yang
belum memiliki produk, pencatatan, atau kapasitas operasional yang
memadai. Sebaliknya, pelaku usaha yang telah menggunakan platform
digital mungkin membutuhkan dukungan lanjutan berupa penguatan strategi
penjualan, kualitas produk, dan perluasan pasar.

Oleh karena itu, diperlukan pendekatan kebijakan yang mempertimbangkan
dua dimensi sekaligus: karakteristik dan kebutuhan pelaku usaha serta
distribusi spasial UMKM antarkelurahan. Pemetaan geospasial dapat
membantu pemerintah mengidentifikasi wilayah yang perlu memperoleh
perhatian lebih lanjut dan menyusun program yang lebih terarah.

\chapter*{Temuan Utama Penelitian}\label{temuan-utama-penelitian}

\section*{Kebutuhan UMKM perlu dibaca secara
multidimensional}\label{kebutuhan-umkm-perlu-dibaca-secara-multidimensional}

Penelitian mengembangkan Indeks Kesiapan Digital (\emph{Digital
Readiness Index} atau DRI) melalui penggabungan dimensi yang mencakup
karakteristik usaha, keterampilan digital, strategi, keamanan, praktik
pengelolaan data, pemanfaatan ekonomi teknologi digital, dan
infrastruktur.

Pendekatan indeks komposit digunakan untuk merangkum beberapa dimensi
kesiapan digital tanpa mengasumsikan bahwa seluruh dimensi tersebut
harus membentuk satu faktor laten hierarkis. Dengan demikian, DRI
berfungsi sebagai ukuran ringkas untuk analisis dan pemetaan, bukan
sebagai pengganti informasi mengenai dimensi penyusunnya.

Implikasi kebijakannya adalah bahwa nilai kesiapan digital perlu dibaca
bersama karakteristik usaha dan indikator penyusunnya. Dua kelurahan
dengan nilai DRI yang relatif serupa dapat memiliki kebutuhan intervensi
yang berbeda apabila hambatan utama UMKM di kedua wilayah tersebut tidak
sama.

\section*{Pemanfaatan perdagangan elektronik berkaitan positif dengan
kesiapan
digital}\label{pemanfaatan-perdagangan-elektronik-berkaitan-positif-dengan-kesiapan-digital}

Dalam model spasial tingkat kelurahan, koefisien pemanfaatan perdagangan
elektronik (\emph{e-commerce}) diperkirakan sebesar 0,396 dengan selang
kredibel 95\% sebesar 0,040--0,754. Seluruh selang berada di atas nol,
sehingga hasil model mendukung adanya hubungan positif antara
pemanfaatan perdagangan elektronik dan DRI, dengan tetap memperhatikan
keterbatasan desain penelitian.

Temuan ini menunjukkan pentingnya mendorong pemanfaatan teknologi
digital dalam kegiatan usaha secara nyata, bukan hanya meningkatkan
pengetahuan umum mengenai teknologi. Pendampingan dapat diarahkan pada
penggunaan platform penjualan, pengelolaan katalog produk, komunikasi
dengan pelanggan, dan pemanfaatan kanal digital yang sesuai dengan
kebutuhan usaha.

Namun, karena analisis bersifat observasional dan menggunakan agregasi
tingkat kelurahan, hasil tersebut tidak membuktikan bahwa penggunaan
perdagangan elektronik secara langsung menyebabkan peningkatan kesiapan
digital. Hubungan yang teramati juga dapat mencerminkan karakteristik
usaha dan faktor lain yang berkaitan.

\section*{Kesiapan halal dan hambatan halal perlu ditangani secara
bersamaan}\label{kesiapan-halal-dan-hambatan-halal-perlu-ditangani-secara-bersamaan}

Model menghasilkan estimasi koefisien positif untuk kesiapan halal
sebesar 0,737 dan koefisien negatif untuk hambatan halal sebesar −0,420.
Kedua estimasi tersebut menunjukkan arah hubungan yang berbeda dan
relevan bagi perumusan intervensi.

Kesiapan halal dapat menjadi pintu masuk untuk mengintegrasikan
pengembangan usaha dengan pendampingan pemenuhan persyaratan halal.
Sementara itu, hambatan halal perlu dipetakan lebih lanjut untuk
mengetahui kendala yang dapat ditangani melalui penyederhanaan akses
informasi, fasilitasi dokumen, pendampingan proses sertifikasi, atau
penguatan praktik produksi.

Besaran koefisien tersebut tidak boleh ditafsirkan sebagai ukuran dampak
kebijakan. Ketidakpastian estimasi dan kekuatan bukti untuk
masing-masing prediktor perlu diperiksa sebelum menetapkan besaran
prioritas intervensi. Secara khusus, hasil model belum cukup untuk
menyatakan bahwa semua aspek kesiapan maupun hambatan halal memiliki
hubungan yang sama kuat dengan DRI.

\section*{Bukti mengenai infrastruktur digital masih memerlukan
pendalaman}\label{bukti-mengenai-infrastruktur-digital-masih-memerlukan-pendalaman}

Koefisien infrastruktur digital diperkirakan sebesar 0,497, dengan
selang kredibel 95\% sebesar −0,164 hingga 1,166. Walaupun arah
estimasinya positif, selang tersebut mencakup nol.

Dengan demikian, penelitian belum memberikan bukti yang cukup kuat untuk
menyimpulkan adanya hubungan statistik yang jelas antara indikator
infrastruktur yang digunakan dan DRI pada tingkat kelurahan.

Temuan ini bukan alasan untuk mengabaikan infrastruktur digital.
Sebaliknya, pemerintah perlu memeriksa apakah indikator yang digunakan
telah menggambarkan kualitas konektivitas, keterjangkauan biaya,
ketersediaan perangkat, dan akses aktual yang dialami pelaku usaha.
Evaluasi tersebut penting agar kebijakan infrastruktur didasarkan pada
hambatan yang benar-benar teridentifikasi.

\section*{Pemetaan geospasial membuka peluang untuk intervensi yang
lebih
terarah}\label{pemetaan-geospasial-membuka-peluang-untuk-intervensi-yang-lebih-terarah}

Penelitian mengagregasikan indikator UMKM untuk analisis pada 50
kelurahan di Kota Mataram dan mengembangkan peta interaktif yang memuat
informasi lokasi serta karakteristik usaha. Pemetaan ini dapat digunakan
sebagai perangkat awal untuk membaca distribusi UMKM dan menyusun
prioritas program berbasis wilayah.

Analisis spasial juga mempertimbangkan keterkaitan antarkelurahan.
Pendekatan ini relevan karena karakteristik UMKM di suatu wilayah tidak
selalu berdiri sendiri dan dapat berkaitan dengan kondisi wilayah di
sekitarnya.

Peta perlu digunakan bersama informasi jumlah UMKM, komposisi skala
usaha, kesiapan digital, dan indikator halal. Wilayah dengan nilai
rata-rata DRI rendah, misalnya, perlu diperiksa lebih lanjut untuk
menentukan apakah intervensi utama yang dibutuhkan berupa pelatihan
dasar, akses infrastruktur, pendampingan halal, atau kombinasi beberapa
program.

Peta tersebut merupakan instrumen pendukung pengambilan keputusan, bukan
dasar tunggal untuk menetapkan penerima bantuan atau menyimpulkan bahwa
seluruh UMKM dalam satu kelurahan memiliki kondisi yang sama.

\chapter*{Rekomendasi Kebijakan}\label{rekomendasi-kebijakan}

\section*{Terapkan zonasi intervensi berdasarkan kebutuhan
UMKM}\label{terapkan-zonasi-intervensi-berdasarkan-kebutuhan-umkm}

Pemerintah Kota Mataram perlu menggunakan peta geospasial dan indikator
penelitian untuk menyusun prioritas intervensi pada tingkat kelurahan.
Penetapan prioritas tidak cukup didasarkan pada nilai DRI semata, tetapi
juga perlu mempertimbangkan jumlah UMKM, skala usaha, pemanfaatan
perdagangan elektronik, kesiapan halal, hambatan halal, serta
karakteristik lokal.

Zonasi sebaiknya menggunakan pendekatan multidimensional. Sebagai
kerangka operasional awal, pemerintah dapat membedakan kebutuhan
intervensi sebagai berikut:

\begin{itemize}
\tightlist
\item
  \textbf{Penguatan dasar digital:} untuk wilayah atau kelompok usaha
  yang memerlukan peningkatan keterampilan digital dasar dan pengenalan
  pemanfaatan teknologi.
\item
  \textbf{Akselerasi perdagangan elektronik:} untuk UMKM yang telah
  memiliki dasar digital dan memerlukan pendampingan dalam penggunaan
  platform penjualan serta pengembangan pasar.
\item
  \textbf{Fasilitasi kesiapan halal:} untuk UMKM yang membutuhkan
  informasi, pemeriksaan kesiapan dokumen, pendampingan proses
  sertifikasi, dan penguatan praktik produksi.
\item
  \textbf{Penguatan terpadu:} untuk kelompok usaha yang membutuhkan
  pendampingan digital dan halal secara bersamaan.
\end{itemize}

Kategori tersebut merupakan kerangka usulan program, bukan klasifikasi
empiris kelurahan yang telah ditetapkan oleh penelitian. Penetapan
kelurahan prioritas harus dilakukan setelah indikator aktual diperiksa
dan diverifikasi bersama perangkat daerah serta pendamping UMKM.

\section*{Alihkan fokus digitalisasi dari pelatihan umum ke pendampingan
penggunaan}\label{alihkan-fokus-digitalisasi-dari-pelatihan-umum-ke-pendampingan-penggunaan}

Program digitalisasi sebaiknya tidak berhenti pada sosialisasi atau
pelatihan satu kali. Pemerintah perlu mengembangkan pendampingan
bertahap yang berorientasi pada penerapan teknologi dalam kegiatan
usaha.

Kegiatan yang dapat diprioritaskan meliputi:

\begin{enumerate}
\def\labelenumi{\arabic{enumi}.}
\tightlist
\item
  Pembuatan dan perbaikan katalog produk digital.
\item
  Penggunaan media sosial dan platform perdagangan elektronik yang
  sesuai dengan jenis usaha.
\item
  Peningkatan keterampilan komunikasi dan pelayanan pelanggan secara
  digital.
\item
  Pengelolaan pesanan, informasi produk, dan pencatatan transaksi.
\item
  Pendampingan lanjutan untuk menilai penggunaan kanal digital setelah
  pelatihan.
\end{enumerate}

Pelaksanaan dapat melibatkan Dinas Koperasi dan UKM, PLUT, perguruan
tinggi, komunitas usaha, serta mitra platform digital. Keberhasilan
program perlu dinilai melalui perubahan praktik usaha, bukan hanya
jumlah peserta yang mengikuti pelatihan.

\section*{Integrasikan pendampingan halal dengan penguatan mutu
produk}\label{integrasikan-pendampingan-halal-dengan-penguatan-mutu-produk}

Layanan halal perlu dikembangkan sebagai bagian dari pembinaan usaha,
bukan sebagai kegiatan administratif yang terpisah dari pengembangan
produk.

Pemerintah dan mitra pendamping dapat mengembangkan layanan terpadu yang
mencakup:

\begin{enumerate}
\def\labelenumi{\arabic{enumi}.}
\tightlist
\item
  Identifikasi status sertifikasi dan kebutuhan masing-masing UMKM.
\item
  Pemeriksaan awal kelengkapan dokumen dan persyaratan yang relevan.
\item
  Informasi mengenai jalur sertifikasi yang sesuai dengan jenis dan
  kondisi usaha.
\item
  Pendampingan proses pengajuan melalui jalur dan lembaga berwenang.
\item
  Penguatan praktik produksi, kebersihan, konsistensi bahan, dan
  dokumentasi yang mendukung penjaminan produk halal.
\item
  Pemantauan perkembangan pengajuan sampai pada status penyelesaian yang
  dapat diverifikasi.
\end{enumerate}

Skema layanan perlu disesuaikan dengan ketentuan yang berlaku, termasuk
persyaratan program sertifikasi yang relevan. Tidak semua pelaku usaha
dapat menggunakan jalur atau skema yang sama.

Masukan dari kegiatan diskusi pemangku kepentingan dalam penelitian juga
memperlihatkan relevansi kebutuhan informasi mengenai proses sertifikasi
halal dan penjaminan mutu produk melalui layanan pendampingan seperti
PLUT. Kebutuhan tersebut dapat dijadikan dasar untuk merancang layanan
yang lebih mudah diakses oleh pelaku UMKM.

\section*{Bangun layanan pendampingan terpadu digital dan
halal}\label{bangun-layanan-pendampingan-terpadu-digital-dan-halal}

Program digitalisasi dan halal sebaiknya tidak dijalankan sebagai dua
agenda yang terpisah. UMKM yang membutuhkan kedua jenis dukungan dapat
memperoleh pendampingan terpadu melalui klinik usaha di tingkat
kecamatan atau kelurahan, dengan jadwal dan mekanisme layanan yang
terkoordinasi.

Pendampingan dapat menghubungkan pengembangan katalog digital dengan
perbaikan informasi produk, identitas usaha, pencatatan bahan, dan
dokumentasi proses produksi. Dengan demikian, kegiatan digitalisasi
dapat mendukung pengembangan pasar, sementara penguatan praktik usaha
mendukung kesiapan memenuhi persyaratan halal.

Model layanan terpadu perlu memperhatikan kapasitas masing-masing UMKM.
Usaha yang masih membutuhkan pembenahan dasar sebaiknya tidak langsung
dibebani target adopsi teknologi yang kompleks atau proses administratif
yang belum siap dipenuhi.

\section*{Bangun sistem pemantauan berbasis data
geospasial}\label{bangun-sistem-pemantauan-berbasis-data-geospasial}

Peta interaktif hasil penelitian perlu dikembangkan sebagai salah satu
instrumen pendukung perencanaan dan evaluasi program UMKM. Pemerintah
dapat mengintegrasikan data spasial dengan informasi program,
karakteristik usaha, status pendampingan, dan perkembangan hasil
intervensi.

Sistem pemantauan perlu memperhatikan kualitas data, pembaruan
informasi, perlindungan data pelaku usaha, dan konsistensi definisi
indikator. Data individu yang bersifat sensitif atau mengidentifikasi
pelaku usaha tidak perlu ditampilkan secara terbuka apabila tidak
dibutuhkan untuk tujuan layanan publik.

Pembaruan berkala memungkinkan pemerintah membandingkan kondisi awal
dengan perkembangan setelah program dilaksanakan. Namun, perubahan
indikator tidak boleh langsung dianggap sebagai dampak program tanpa
mempertimbangkan faktor lain yang mungkin memengaruhi hasil.

\chapter*{Rencana Implementasi}\label{rencana-implementasi}

Pelaksanaan rekomendasi dapat dilakukan secara bertahap melalui
koordinasi lintas perangkat daerah dan mitra pendamping.

\textbf{Tahap 1: Validasi data dan penetapan prioritas.} Pemerintah
bersama tim peneliti dan pendamping UMKM meninjau peta, memeriksa
kelengkapan data, mengidentifikasi kebutuhan yang paling mendesak, dan
menyepakati kriteria penetapan wilayah prioritas. Hasilnya berupa daftar
wilayah dan kelompok UMKM yang memerlukan jenis intervensi tertentu.

\textbf{Tahap 2: Perancangan dan pelaksanaan program.} Perangkat daerah
bersama PLUT, perguruan tinggi, dan mitra terkait menyelenggarakan
pendampingan digital, konsultasi halal, serta penguatan mutu produk
sesuai dengan kebutuhan yang telah diidentifikasi. Program dapat dimulai
pada wilayah prioritas yang datanya telah diverifikasi.

\textbf{Tahap 3: Pemantauan dan evaluasi.} Perkembangan peserta dipantau
melalui indikator yang telah ditetapkan. Hasil evaluasi digunakan untuk
memperbaiki materi, mekanisme pendampingan, dan distribusi sumber daya
pada periode berikutnya.

Pembagian peran yang disarankan adalah sebagai berikut:

{\def\LTcaptype{none} % do not increment counter
\begin{longtable}[]{@{}
  >{\raggedright\arraybackslash}p{(\linewidth - 2\tabcolsep) * \real{0.5000}}
  >{\raggedright\arraybackslash}p{(\linewidth - 2\tabcolsep) * \real{0.5000}}@{}}
\toprule\noalign{}
\begin{minipage}[b]{\linewidth}\raggedright
Pemangku kepentingan
\end{minipage} & \begin{minipage}[b]{\linewidth}\raggedright
Peran yang disarankan
\end{minipage} \\
\midrule\noalign{}
\endhead
\bottomrule\noalign{}
\endlastfoot
Pemerintah Kota Mataram dan perangkat daerah terkait & Menetapkan
prioritas, mengoordinasikan program, dan mengintegrasikan hasil pemetaan
ke dalam perencanaan. \\
Dinas Koperasi dan UKM serta PLUT & Menyelenggarakan pendampingan usaha,
fasilitasi digitalisasi, konsultasi, dan rujukan layanan halal. \\
Lembaga dan pendamping halal yang berwenang & Memberikan informasi
persyaratan, mendampingi proses sesuai kewenangan, dan memverifikasi
perkembangan sertifikasi. \\
Perguruan tinggi & Mendukung analisis data, pendampingan teknis,
evaluasi program, dan pembaruan informasi geospasial. \\
Pelaku dan komunitas UMKM & Mengidentifikasi kebutuhan, mengikuti
pendampingan, menerapkan praktik yang relevan, dan memberikan umpan
balik. \\
\end{longtable}
}

\chapter*{Indikator Pemantauan dan
Evaluasi}\label{indikator-pemantauan-dan-evaluasi}

Agar pelaksanaan kebijakan dapat dievaluasi secara terukur, pemerintah
perlu menetapkan indikator awal, target, periode pemantauan, dan sumber
data sebelum program dimulai.

{\def\LTcaptype{none} % do not increment counter
\begin{longtable}[]{@{}
  >{\raggedright\arraybackslash}p{(\linewidth - 2\tabcolsep) * \real{0.5000}}
  >{\raggedright\arraybackslash}p{(\linewidth - 2\tabcolsep) * \real{0.5000}}@{}}
\toprule\noalign{}
\begin{minipage}[b]{\linewidth}\raggedright
Dimensi
\end{minipage} & \begin{minipage}[b]{\linewidth}\raggedright
Indikator yang disarankan
\end{minipage} \\
\midrule\noalign{}
\endhead
\bottomrule\noalign{}
\endlastfoot
Jangkauan intervensi & Jumlah kelurahan yang dipetakan dan divalidasi;
jumlah UMKM yang menerima pendampingan; distribusi peserta menurut skala
usaha. \\
Kesiapan digital & Perubahan DRI dan dimensi penyusunnya; perubahan
keterampilan atau praktik digital yang terukur. \\
Pemanfaatan perdagangan elektronik & Proporsi UMKM yang aktif
menggunakan kanal digital untuk penjualan; keberlanjutan penggunaan
setelah pendampingan. \\
Infrastruktur digital & Kendala konektivitas, perangkat, dan
keterjangkauan yang teridentifikasi; tindak lanjut terhadap hambatan
yang dapat ditangani. \\
Kesiapan halal & Jumlah UMKM yang telah menjalani pemeriksaan kesiapan
dan menerima pendampingan sesuai kebutuhan. \\
Progres sertifikasi halal & Jumlah pengajuan yang dimulai, diproses, dan
diselesaikan; status sertifikasi yang dapat diverifikasi. \\
Mutu produk & Perbaikan dokumentasi dan praktik produksi yang relevan;
tindak lanjut atas hasil pembinaan mutu. \\
Efektivitas layanan & Tingkat penyelesaian pendampingan, waktu layanan,
kepuasan peserta, dan tindak lanjut hasil evaluasi. \\
\end{longtable}
}

Nilai awal indikator perlu diukur sebelum intervensi dilaksanakan.
Target kuantitatif sebaiknya ditetapkan setelah jumlah sasaran,
kapasitas pendamping, anggaran, dan kondisi awal diketahui, bukan
ditentukan secara arbitrer.

Evaluasi dapat dilakukan dengan membandingkan indikator sebelum dan
sesudah program. Apabila sumber daya memungkinkan, rancangan evaluasi
yang memiliki kelompok pembanding atau penerapan bertahap dapat membantu
membedakan perubahan yang berkaitan dengan intervensi dari perubahan
yang terjadi karena faktor lain.

\chapter*{Implikasi dan Batasan
Kebijakan}\label{implikasi-dan-batasan-kebijakan}

Penelitian ini menyediakan dasar empiris untuk mengembangkan kebijakan
UMKM yang lebih terarah secara spasial. Namun, beberapa batasan perlu
dipertimbangkan dalam penerapannya.

Pertama, indeks komposit merangkum beberapa dimensi kesiapan digital.
Nilai agregat perlu dibaca bersama indikator penyusunnya agar perbedaan
kebutuhan antarwilayah tidak terabaikan.

Kedua, hasil pemodelan spasial menunjukkan hubungan statistik pada
tingkat agregasi kelurahan. Hasil tersebut tidak membuktikan hubungan
sebab-akibat pada tingkat individu. Oleh sebab itu, hasil model perlu
dipadukan dengan verifikasi lapangan sebelum digunakan untuk menetapkan
intervensi spesifik.

Ketiga, ketidakpastian estimasi berbeda antarvariabel. Temuan mengenai
hubungan positif perdagangan elektronik memperoleh dukungan statistik
dari selang kredibelnya, sedangkan hasil untuk infrastruktur digital
masih belum konklusif. Tingkat kepastian hasil perlu dipertimbangkan
dalam menentukan prioritas program.

Keempat, peta wilayah tidak menggambarkan seluruh variasi kondisi UMKM
di dalam satu kelurahan. Pengambilan keputusan mengenai peserta program
perlu didukung pemeriksaan kebutuhan pada tingkat usaha.

Kelima, status sertifikasi halal dan kesiapan memenuhi persyaratan halal
merupakan dua hal yang berbeda. Program perlu membedakan pelaku usaha
yang telah memiliki sertifikat, yang sedang menjalani proses, dan yang
masih memerlukan pendampingan awal.

Dengan mempertimbangkan batasan tersebut, hasil penelitian sebaiknya
digunakan sebagai dasar untuk menyusun prioritas awal, melakukan
validasi, dan memperbaiki kebijakan secara berkelanjutan.

\chapter*{Penutup}\label{penutup}

Akselerasi digitalisasi dan sertifikasi halal UMKM Kota Mataram
membutuhkan pendekatan yang menghubungkan data, layanan pendampingan,
dan kebutuhan aktual pelaku usaha. Hasil penelitian menunjukkan bahwa
pemanfaatan perdagangan elektronik dan kesiapan halal merupakan dimensi
yang relevan untuk dipertimbangkan dalam perumusan kebijakan, sementara
hambatan halal memerlukan perhatian khusus dan peran infrastruktur perlu
dievaluasi lebih lanjut berdasarkan kondisi lapangan.

Pemerintah Kota Mataram dapat memanfaatkan pemetaan geospasial sebagai
instrumen untuk menentukan prioritas wilayah, mengoordinasikan layanan,
dan memantau perkembangan program. Pendekatan ini diharapkan membantu
mengarahkan sumber daya secara lebih tepat, memperkuat kapasitas usaha,
serta membangun ekosistem UMKM yang lebih adaptif terhadap perkembangan
ekonomi digital dan ketentuan halal.




\end{document}
